% =====================================================================
%  defs.tex — Macros del banc de preguntes
%  ─────────────────────────────────────────────────────────────────────
%  Va sempre després de headers.tex i abans de \begin{document}, amb
%  \newif\ifsolucions ja declarat. Vegeu headers.tex.
% =====================================================================

% ── 0. Segell de versió ──────────────────────────────────────────────
% build.py hi afegeix, al final, \def\bancversio{...}. Cada examen que
% baixa el lloc comença amb \bancrequereix{...}: si el defs.tex és d'una
% altra versió, LaTeX avisa i cal tornar a baixar l'entorn.
\makeatletter
\newcommand{\bancrequereix}[1]{%
  \def\bp@v{#1}%
  \ifx\bp@v\bancversio\else
    \GenericWarning{}{Banc de preguntes: aquest defs.tex (\bancversio) no és el
      que va generar l'examen (#1). Torna a baixar headers.tex i defs.tex}%
  \fi}
\makeatother

% ── 0b. Capçalera d'examen: buida al banc ────────────────────────────
% El banc no porta cap dada de cap centre: ni logo, ni segell, ni curs, ni
% departament. Això viu NOMÉS a la carpeta d'exàmens del professorat, en un
% capsalera.tex que main.tex incorpora si hi és. Aquí, la macro no escriu
% res, i per això els PDF del banc i el lloc web no ensenyen cap capçalera.
\providecommand{\capsaleraexamen}{}

% ── 1. Capçalera de pregunta ──────────────────────────────────────────
% El filet va A SOBRE de cada pregunta i el \Needspace el manté enganxat
% al seu text: així el filet mai queda orfe al capdamunt d'una pàgina.
\newcommand{\encapcalament}[1]{%
  \par\Needspace{8\baselineskip}%
  \bigskip\hrule\medskip
  \noindent\textbf{#1}\par\smallskip}

% ── 1b. Ela geminada (l·l) ───────────────────────────────────────────
% El punt volat es compon com un símbol solt i queda massa separat de les
% eles: «paral · lela». Amb una mica de kerning, «paral·lela».
\DeclareUnicodeCharacter{00B7}{\kern-0.1em\textperiodcentered\kern-0.1em}

% ── 2. Apartats a) b) c) amb la seva puntuació ────────────────────────
% \apartat{0,75} escriu "(0,75 punts)".  build.py llegeix aquests valors
% del .tex: la puntuació viu AQUÍ i enlloc més. A les fonts, \apartat[1,25]{0,75}
% porta també la puntuació de 50 min, però materialitza() el deixa com a
% \apartat{0,75} o \apartat{1,25} abans d'arribar aquí.
\makeatletter
\newcommand{\bp@ptsmot}[1]{\def\bp@a{#1}\def\bp@u{1}%
  \ifx\bp@a\bp@u 1~punt\else #1~punts\fi}
\newcommand{\apartat}[1]{\item \textit{(\bp@ptsmot{#1})}\enspace\ignorespaces}
\makeatother
\newenvironment{apartats}
  {\begin{enumerate}[label=\textbf{\alph*)},leftmargin=*,itemsep=\bigskipamount,
                     topsep=\medskipamount,parsep=\smallskipamount]}
  {\end{enumerate}}

% ── 3. Graella de subapartats i) ii) iii) en mode display ─────────────
% \begin{graella}{4} \sa ... & \sa ... \\ \sa ... \end{graella}
% La columna força \displaystyle: els \lim hi surten amb el subíndex A
% SOTA, igual que en una fórmula destacada. Mai fem servir array pelat.
\newcounter{bpsa}
\newcommand{\sa}{\stepcounter{bpsa}\textup{\roman{bpsa})}~}
\newenvironment{graella}[1]
  {\setcounter{bpsa}{0}\setlength{\arraycolsep}{1.2em}%
   \[\begin{array}{*{#1}{>{\displaystyle}l}}}
  {\end{array}\]}

% ── 3b. Fórmules destacades: sempre el mateix aire ───────────────────
% Si la línia d'abans és curta («determina:»), TeX fa servir un espai
% «curt», gairebé zero, i la fórmula queda enganxada al text. L'igualem
% a l'espai normal.
\makeatletter
\g@addto@macro\normalsize{%
  \setlength\abovedisplayshortskip{\abovedisplayskip}%
  \setlength\belowdisplayshortskip{\belowdisplayskip}}
\makeatother
\AtBeginDocument{\normalsize}

% ── 3c. Color de les gràfiques ───────────────────────────────────────
% Totes les gràfiques del banc dibuixen les funcions amb aquest color, i
% cap pregunta no n'escriu cap directament. Es pot canviar des de la
% carpeta d'exàmens: \renewcommand{\colorgrafica}{red} al capsalera.tex.
\newcommand{\colorgrafica}{blue}

% ── 4. Condicions dins de \begin{cases} ───────────────────────────────
% \si{-1\le x\le 2}  →  el signe menys surt unari i ben espaiat.
% Escriure \text{si } -1\le x\le 2 a pèl produeix "si − 1 ≤ x ≤ 2". Bug.
\newcommand{\si}[1]{\text{si }{#1}}

% ── 5. Solucions ──────────────────────────────────────────────────────
% REGLA: \end{solucio} ha d'anar SOL a la seva línia (ho exigeix el
% paquet comment quan la solució s'exclou). build.py ho comprova.
\ifsolucions
  \newenvironment{solucio}
    {\par\nopagebreak\medskip\begingroup\color{blue!55!black}\small
     \textbf{Solució.}\enspace\ignorespaces}
    {\par\endgroup\smallskip}
\else
  \excludecomment{solucio}
\fi

% ── 6. Procedència de les preguntes PAU ───────────────────────────────
% La línia «PAU juny 2026, sèrie 1» la injecta el build (i el lloc) just
% després de la capçalera, a partir del codi de la pregunta i de
% pau/convocatories.json. No s'escriu mai a mà: build.py ho rebutja.
\newcommand{\procedencia}[1]{\noindent{\small\itshape #1}\par\vspace{3pt}}
